% !TEX program = xelatex
% Préambule de l'édition du Livre-Journal d'Hippolyte Combes (généré par edition/src/generer.py)
\documentclass[10pt,twoside,openany]{book}
\usepackage[paperwidth=160mm,paperheight=226mm,inner=10.5mm,outer=8.5mm,top=7.5mm,bottom=8.5mm,
            includehead,headheight=10pt,headsep=2.6mm,nofoot]{geometry}
\usepackage{fontspec}
\setmainfont{LibertinusSerif}[
  Extension = .otf,
  UprightFont = *-Regular,
  ItalicFont = *-Italic,
  BoldFont = *-Bold,
  BoldItalicFont = *-BoldItalic,
  Ligatures = {Common, Historic, TeX},
  Numbers = OldStyle,
  SmallCapsFeatures = {Letters=SmallCaps}
]
\setsansfont{LibertinusSans}[
  Extension = .otf,
  UprightFont = *-Regular,
  ItalicFont = *-Italic,
  BoldFont = *-Bold
]
% Transcription diplomatique : sans ligatures historiques, chiffres alignés
\newfontfamily\LJpolice{LibertinusSerif}[
  Extension = .otf,
  UprightFont = *-Regular,
  ItalicFont = *-Italic,
  BoldFont = *-Bold,
  BoldItalicFont = *-BoldItalic,
  Ligatures = {Common, TeX},
  Numbers = Lining
]
\newfontfamily\LJchiffres{LibertinusSerif}[
  Extension = .otf, UprightFont = *-Regular, BoldFont = *-Bold, Numbers = Lining]
\usepackage[french]{babel}
\frenchsetup{StandardItemLabels=true}
\usepackage{microtype}
\usepackage{graphicx}
\usepackage[table]{xcolor}
\usepackage{tikz}
\usetikzlibrary{decorations.pathreplacing,calc,positioning,arrows.meta}
\usepackage{fancyhdr}
\usepackage{titlesec}
\usepackage{array,booktabs,longtable,multicol}
\usepackage{enumitem}
\usepackage{xparse}
\usepackage{yfonts}
\usepackage[hyphens]{url}
\usepackage{metalogo}
\usepackage{hyperref}
\hypersetup{colorlinks=true,linkcolor=black,citecolor=black,urlcolor=black!70!red,
  pdfauthor={Hippolyte Combes (1876–1970)},
  pdftitle={Hippolyte Combes, Livre-Journal (25 juillet 1916 – 30 avril 1917). Édition avec fac-similés, commentaire critique et historique, et illustrations}}

\definecolor{LJgris}{gray}{0.35}
\definecolor{LJfilet}{gray}{0.55}
\definecolor{LJencre}{gray}{0.05}

\setlength{\parindent}{1.1em}
\setlength{\parskip}{0pt}
\setlength{\emergencystretch}{2em}
\frenchspacing
\widowpenalty=10000 \clubpenalty=10000

% ---------------------------------------------------------------------------
% Titres
% ---------------------------------------------------------------------------
\titleformat{\chapter}[display]{\normalfont\centering}
  {\small\scshape\MakeLowercase{\chaptertitlename}~\Roman{chapter}}{2mm}
  {\Large\scshape}[\vspace{1mm}\rule{18mm}{0.4pt}]
\titlespacing*{\chapter}{0pt}{0mm}{6mm}
\titleformat{\section}{\normalfont\large\scshape}{\thesection.}{0.6em}{}
\titleformat{\subsection}{\normalfont\itshape}{}{0em}{}
\titlespacing*{\section}{0pt}{4mm}{2mm}
\titlespacing*{\subsection}{0pt}{3mm}{1mm}
\renewcommand{\thechapter}{\Roman{chapter}}
\renewcommand{\thesection}{\Roman{section}}
\renewcommand{\thefigure}{\arabic{figure}}
\makeatletter
\renewcommand*\l@chapter[2]{%
  \ifnum \c@tocdepth >\m@ne
    \addpenalty{-\@highpenalty}%
    \vskip 1.0em \@plus\p@
    \setlength\@tempdima{2.2em}%
    \begingroup
      \parindent \z@ \rightskip \@pnumwidth
      \parfillskip -\@pnumwidth
      \leavevmode \bfseries
      \advance\leftskip\@tempdima
      \hskip -\leftskip
      #1\nobreak\hfil
      \nobreak\hb@xt@\@pnumwidth{\hss #2\kern-\p@\kern\p@}\par
      \penalty\@highpenalty
    \endgroup
  \fi}
\makeatother
\setcounter{secnumdepth}{1}

% ---------------------------------------------------------------------------
% En-têtes courants
% ---------------------------------------------------------------------------
\newcommand{\LJmois}{}
\newcommand{\LJfolio}{}
\newcommand{\LJtitrecourant}{\textsc{Hippolyte Combes} — \textit{Livre-Journal}}
\pagestyle{fancy}
\fancyhf{}
\fancyhead[LE]{\footnotesize\LJchiffres\thepage\hfill\LJtitrecourant\hfill\LJmois}
\fancyhead[RO]{\footnotesize\LJmois\hfill\LJfolio\hfill\LJchiffres\thepage}
\renewcommand{\headrulewidth}{0.3pt}
\fancypagestyle{plain}{\fancyhf{}\fancyhead[LE,RO]{\footnotesize\LJchiffres\thepage}\renewcommand{\headrulewidth}{0pt}}
\fancypagestyle{vide}{\fancyhf{}\renewcommand{\headrulewidth}{0pt}}

% saut vers une page paire (gauche), pour garantir le vis-à-vis
\newcommand{\LJpagepaire}{\clearpage\ifodd\value{page}\hbox{}\thispagestyle{vide}\newpage\fi}
\newcommand{\LJpageimpaire}{\clearpage\ifodd\value{page}\else\hbox{}\thispagestyle{vide}\newpage\fi}

% ---------------------------------------------------------------------------
% Signes diplomatiques (en ligne)
% ---------------------------------------------------------------------------
\newlength{\LJsoul}
\newcommand{\LJsl}[1]{\leavevmode\setbox0\hbox{#1}\rlap{\rule[-0.32ex]{\wd0}{0.42pt}}\box0\relax}
\newcommand{\LJsll}[1]{\leavevmode\setbox0\hbox{#1}\rlap{\rule[-0.32ex]{\wd0}{0.42pt}}\rlap{\rule[-0.62ex]{\wd0}{0.42pt}}\box0\relax}
\newcommand{\LJsup}[1]{\textsuperscript{#1}}
\newcommand{\LJdel}[1]{\leavevmode\setbox0\hbox{#1}\rlap{\rule[0.55ex]{\wd0}{0.5pt}}\box0\relax}
\newcommand{\LJsur}[1]{#1}
\newcommand{\LJadd}[1]{\textsuperscript{\textbackslash{}#1/}}
\newcommand{\LJaddm}[1]{\textsuperscript{\textbackslash{}#1/}}
\newcommand{\LJinc}[1]{#1\kern0.06em\textsuperscript{\textup{\color{LJgris}?}}}
\NewDocumentCommand{\LJill}{m}{[\textit{\ifx\relax#1\relax illis.\else #1~mots illis.\fi}]}
\newcommand{\LJlac}{[\textit{lacune}]}
\newcommand{\LJtrait}{\leaders\hbox{\rule[0.48ex]{1pt}{0.42pt}}\hfill\kern0.2em}
\newcommand{\LJfrac}[2]{\textsuperscript{#1}\kern-0.05em/\kern-0.05em\textsubscript{#2}}

% ---------------------------------------------------------------------------
% Double page : transcription (page paire) / fac-similé (page impaire)
% ---------------------------------------------------------------------------
\newlength{\LJcol}       \setlength{\LJcol}{88mm}   % colonne de transcription
\newlength{\LJgap}       \setlength{\LJgap}{4mm}  % gouttière
\newlength{\LJappw}                                 % colonne d'apparat
\newlength{\LJpitch}                                % interligne de la transcription
\newlength{\LJappsep}    \setlength{\LJappsep}{0.9pt}
\newdimen\LJappcur
\newbox\LJboxT \newbox\LJboxA \newbox\LJboxP \newbox\LJboxN
\newcount\LJessai
\newcommand{\LJtaille}{}
\newcommand{\LJlabelcourant}{}
\newcommand{\LJpoliceapp}{}
\newcommand{\LJpolicereg}{}

% Jeux de tailles, du plus grand au plus petit : <transcription>/<interligne> ; apparat ; registres
\newcommand{\LJjeu}[1]{%
  \ifcase#1\or
    \renewcommand{\LJtaille}{\fontsize{11pt}{13.8pt}\selectfont}\setlength{\LJpitch}{13.8pt}%
    \renewcommand{\LJpoliceapp}{\fontsize{8.2pt}{9.6pt}\selectfont}%
    \renewcommand{\LJpolicereg}{\fontsize{8.6pt}{10.4pt}\selectfont}%
  \or
    \renewcommand{\LJtaille}{\fontsize{10.5pt}{13.1pt}\selectfont}\setlength{\LJpitch}{13.1pt}%
    \renewcommand{\LJpoliceapp}{\fontsize{8pt}{9.3pt}\selectfont}%
    \renewcommand{\LJpolicereg}{\fontsize{8.4pt}{10.1pt}\selectfont}%
  \or
    \renewcommand{\LJtaille}{\fontsize{10pt}{12.5pt}\selectfont}\setlength{\LJpitch}{12.5pt}%
    \renewcommand{\LJpoliceapp}{\fontsize{7.8pt}{9.1pt}\selectfont}%
    \renewcommand{\LJpolicereg}{\fontsize{8.2pt}{9.8pt}\selectfont}%
  \or
    \renewcommand{\LJtaille}{\fontsize{9.5pt}{11.9pt}\selectfont}\setlength{\LJpitch}{11.9pt}%
    \renewcommand{\LJpoliceapp}{\fontsize{7.6pt}{8.8pt}\selectfont}%
    \renewcommand{\LJpolicereg}{\fontsize{8pt}{9.5pt}\selectfont}%
  \or
    \renewcommand{\LJtaille}{\fontsize{9pt}{11.2pt}\selectfont}\setlength{\LJpitch}{11.2pt}%
    \renewcommand{\LJpoliceapp}{\fontsize{7.3pt}{8.5pt}\selectfont}%
    \renewcommand{\LJpolicereg}{\fontsize{7.7pt}{9.2pt}\selectfont}%
  \or
    \renewcommand{\LJtaille}{\fontsize{8.6pt}{10.6pt}\selectfont}\setlength{\LJpitch}{10.6pt}%
    \renewcommand{\LJpoliceapp}{\fontsize{7.1pt}{8.2pt}\selectfont}%
    \renewcommand{\LJpolicereg}{\fontsize{7.4pt}{8.8pt}\selectfont}%
  \or
    \renewcommand{\LJtaille}{\fontsize{8.2pt}{10pt}\selectfont}\setlength{\LJpitch}{10pt}%
    \renewcommand{\LJpoliceapp}{\fontsize{6.9pt}{7.9pt}\selectfont}%
    \renewcommand{\LJpolicereg}{\fontsize{7.1pt}{8.4pt}\selectfont}%
  \or
    \renewcommand{\LJtaille}{\fontsize{7.8pt}{9.5pt}\selectfont}\setlength{\LJpitch}{9.5pt}%
    \renewcommand{\LJpoliceapp}{\fontsize{6.7pt}{7.6pt}\selectfont}%
    \renewcommand{\LJpolicereg}{\fontsize{6.9pt}{8.1pt}\selectfont}%
  \or
    \renewcommand{\LJtaille}{\fontsize{7.4pt}{9pt}\selectfont}\setlength{\LJpitch}{9pt}%
    \renewcommand{\LJpoliceapp}{\fontsize{6.5pt}{7.3pt}\selectfont}%
    \renewcommand{\LJpolicereg}{\fontsize{6.7pt}{7.8pt}\selectfont}%
  \else\LJjeu{9}\fi}
\newcommand{\LJnbjeux}{9}

% Rangées de la colonne de transcription. Hauteur d'une ligne de texte = \LJpitch (contenu « écrasé »
% pour que la géométrie reste exacte : la position verticale de chaque ligne est connue du générateur).
% #1 numéro (vide si non imprimé) ; #2 retrait (em) ; #3 alignement g/c/d ; #4 signes (TikZ) ; #5 texte
\newcommand{\LJnum}[1]{\llap{\LJchiffres\fontsize{6pt}{6pt}\selectfont\color{LJgris}#1\hspace{5.2mm}}}
\newcommand{\LJt}[5]{%
  \setbox\LJboxN\hbox{\LJtaille\LJpolice\NoAutoSpacing\hskip#2em\relax#5}%
  \ifdim\wd\LJboxN>\LJcol \setbox\LJboxN\hbox{\resizebox{\LJcol}{!}{\box\LJboxN}}\fi
  \hbox to \LJcol{\vrule height 0.76\LJpitch depth 0.24\LJpitch width 0pt\LJnum{#1}#4%
    \if#3g\else\hfil\fi\smash{\box\LJboxN}\if#3d\else\hfil\fi}}
% Ligne contenant \hfill : la boîte occupe toute la largeur
\newcommand{\LJtp}[5]{%
  \setbox\LJboxN\hbox{\LJtaille\LJpolice\NoAutoSpacing\hskip#2em\relax#5}%
  \ifdim\wd\LJboxN>\LJcol
    \setbox\LJboxN\hbox{\resizebox{\LJcol}{!}{\box\LJboxN}}%
  \else
    \setbox\LJboxN\hbox to \LJcol{\LJtaille\LJpolice\NoAutoSpacing\hskip#2em\relax#5}%
  \fi
  \hbox to \LJcol{\vrule height 0.76\LJpitch depth 0.24\LJpitch width 0pt\LJnum{#1}#4\smash{\box\LJboxN}\hss}}
% Filet : #1 position g/c/d ; #2 fraction de la largeur
\newcommand{\LJf}[2]{%
  \hbox to \LJcol{\vrule height 0.28\LJpitch depth 0.22\LJpitch width 0pt
    \if#1c\hfil\fi\if#1d\hfil\fi
    \color{LJencre}\rule[0.04\LJpitch]{#2\LJcol}{0.5pt}%
    \if#1c\hfil\fi\if#1g\hfil\fi}}
% Blanc : #1 = nombre de lignes
\newcommand{\LJb}[1]{\hbox to \LJcol{\vrule height #1\LJpitch depth 0pt width 0pt\hfil}}

% Signes couvrant plusieurs lignes (dessinés en surimpression) ; #1 = hauteur totale en interlignes
\newcommand{\LJcrochet}[1]{\rlap{\begin{tikzpicture}[overlay,line width=0.5pt,color=LJencre]
  \draw (-0.6mm,0.74\LJpitch) -- ++(-0.9mm,0) -- ($(-1.5mm,-#1\LJpitch+0.74\LJpitch+0.06\LJpitch)$) -- ++(0.9mm,0);
  \end{tikzpicture}}}
\newcommand{\LJcrochetd}[1]{\rlap{\hskip\LJcol\begin{tikzpicture}[overlay,line width=0.5pt,color=LJencre]
  \draw (0.6mm,0.74\LJpitch) -- ++(0.9mm,0) -- ($(1.5mm,-#1\LJpitch+0.74\LJpitch+0.06\LJpitch)$) -- ++(-0.9mm,0);
  \end{tikzpicture}}}
\newcommand{\LJbarre}[1]{\rlap{\begin{tikzpicture}[overlay,line width=0.5pt,color=LJencre]
  \draw (-1.2mm,0.74\LJpitch) -- ($(-1.2mm,-#1\LJpitch+0.74\LJpitch+0.06\LJpitch)$);
  \end{tikzpicture}}}
\newcommand{\LJequerre}[1]{\rlap{\begin{tikzpicture}[overlay,line width=0.5pt,color=LJencre]
  \draw (0.2mm,0.74\LJpitch) -- (-1.4mm,0.74\LJpitch) -- ($(-1.4mm,-#1\LJpitch+0.74\LJpitch+0.06\LJpitch)$);
  \end{tikzpicture}}}
\newcommand{\LJequerred}[1]{\rlap{\hskip\LJcol\begin{tikzpicture}[overlay,line width=0.5pt,color=LJencre]
  \draw (-0.2mm,0.74\LJpitch) -- (1.4mm,0.74\LJpitch) -- ($(1.4mm,-#1\LJpitch+0.74\LJpitch+0.06\LJpitch)$);
  \end{tikzpicture}}}
\newcommand{\LJaccoladeg}[1]{\rlap{\begin{tikzpicture}[overlay,line width=0.45pt,color=LJencre]
  \draw[decorate,decoration={brace,amplitude=1.1mm}] ($(-0.7mm,-#1\LJpitch+0.74\LJpitch+0.06\LJpitch)$) -- (-0.7mm,0.74\LJpitch);
  \end{tikzpicture}}}
\NewDocumentCommand{\LJaccolade}{mm}{\rlap{\hskip\LJcol\begin{tikzpicture}[overlay,line width=0.45pt,color=LJencre]
  \draw[decorate,decoration={brace,amplitude=1.1mm}] (0.7mm,0.74\LJpitch) -- ($(0.7mm,-#1\LJpitch+0.74\LJpitch+0.06\LJpitch)$)
    node[midway,right=1.4mm,inner sep=0pt,font=\LJtaille\LJpolice,color=black]{#2};
  \end{tikzpicture}}}

% Apparat marginal : chaque note est placée à la hauteur de sa ligne (#1, en interlignes, ligne de base),
% ou plus bas si la note précédente occupe encore la place.
\newcommand{\LJapp}[2]{%
  \setbox\LJboxN\vtop{\hsize\LJappw\LJpoliceapp\raggedright\hyphenpenalty=600
    \leftskip=1.2em\parindent=-1.2em\strut#2\strut\par}%
  \dimen2=#1\LJpitch
  \dimen4=\dimexpr\LJappcur+\ht\LJboxN\relax
  \ifdim\dimen2<\dimen4 \dimen2=\dimen4 \fi
  \vskip\dimexpr\dimen2-\ht\LJboxN-\LJappcur\relax
  \dimen6=\dp\LJboxN
  \box\LJboxN
  \global\LJappcur=\dimexpr\dimen2+\dimen6+\LJappsep\relax
  \vskip\LJappsep}

% Corps d'une page de transcription ; #1 en-tête ; #2 rangées ; #3 notes d'apparat ; #4 registres
\newcommand{\LJcorps}[4]{%
  \setlength{\LJappw}{\dimexpr\textwidth-\LJcol-\LJgap\relax}%
  #1\par
  \setbox\LJboxT\vtop{\offinterlineskip\hrule height 0pt #2}%
  \global\LJappcur=0pt
  \setbox\LJboxA\vtop{\offinterlineskip\hrule height 0pt #3}%
  \noindent\hbox{\box\LJboxT\hskip\LJgap\box\LJboxA}\par
  \vspace{2.2mm plus 1fill}
  {\LJpolicereg #4}\par}

\newif\ifLJtrop
\newif\ifLJencore
\newlength{\LJpitchbase}
\newcount\LJdil
\newcommand{\LJpagetranscription}[4]{%
  \LJessai=0
  \loop
    \advance\LJessai by 1
    \LJjeu{\the\LJessai}%
    \setbox\LJboxP\vbox{\hsize\textwidth\LJcorps{#1}{#2}{#3}{#4}}%
    \LJtropfalse
    \ifdim\dimexpr\ht\LJboxP+\dp\LJboxP\relax>\textheight \ifnum\LJessai<\LJnbjeux\relax \LJtroptrue\fi\fi
  \ifLJtrop\repeat
  \typeout{LJJEU \LJlabelcourant\space \the\LJessai\space \the\dimexpr\ht\LJboxP+\dp\LJboxP\relax\space \the\textheight}%
  % Dilatation de l'interligne (par pas de 4 %, jusqu'à +28 %) pour remplir la page sans changer de corps
  \setlength{\LJpitchbase}{\LJpitch}\LJdil=100 \LJencoretrue
  \loop
    \advance\LJdil by 4
    \setlength{\LJpitch}{\dimexpr\LJpitchbase*\LJdil/100\relax}%
    \setbox\LJboxP\vbox{\hsize\textwidth\LJcorps{#1}{#2}{#3}{#4}}%
    \ifdim\dimexpr\ht\LJboxP+\dp\LJboxP\relax>\textheight
      \advance\LJdil by -4 \LJencorefalse
    \else\ifnum\LJdil>124 \LJencorefalse\fi\fi
  \ifLJencore\repeat
  \setlength{\LJpitch}{\dimexpr\LJpitchbase*\LJdil/100\relax}%
  \typeout{LJDIL \LJlabelcourant\space \the\LJdil}%
  \vbox to \textheight{\hsize\textwidth\LJcorps{#1}{#2}{#3}{#4}}%
  \newpage}

% En-tête de la page de transcription
% #1 feuillet (ex. « f. 18 r° ») ; #2 cliché (non affiché) ; #3 date ; #4 lieux ; #5 sommaire
\newcommand{\LJentete}[5]{%
  \noindent{\small\scshape #1}\quad{\small #3}\par
  \ifx\relax#4\relax\else\noindent{\footnotesize\itshape #4}\par\fi
  \ifx\relax#5\relax\else\noindent{\footnotesize\color{LJgris}#5}\par\fi
  \vspace{0.6mm}\noindent{\color{LJgris}\rule{\textwidth}{0.35pt}}\par\vspace{1.4mm}}

% Registres du bas de page (sans intitulés, séparés par des filets)
\newcommand{\LJdebutreg}{\noindent{\color{LJgris}\rule{\textwidth}{0.35pt}}\par\vspace{0.7mm}}
\newcommand{\LJsepreg}[1]{\par\vspace{0.6mm}\noindent{\color{LJgris}\rule{#1}{0.3pt}}\par\vspace{0.5mm}}
\newcommand{\LJregcritique}[1]{\noindent #1\par}
\newcommand{\LJregloci}[1]{\noindent #1\par}
\newcommand{\LJreghistoire}[1]{{\leftskip=4mm\noindent #1\par}}

% Page de fac-similé (pleine page ; la légende figure dans l'en-tête courant) ; #1 fichier image
\newcommand{\LJfacsimile}[1]{%
  \vbox to \textheight{\hsize\textwidth\vfil
    \centerline{\includegraphics[width=\textwidth,height=\textheight,keepaspectratio]{#1}}%
    \vfil}%
  \newpage}
% Fac-similé d'une page écrite en travers : image pivotée pour occuper la page (lire en tournant le livre)
\newcommand{\LJfacsimilepaysage}[1]{%
  \vbox to \textheight{\hsize\textwidth\vfil
    \centerline{\includegraphics[angle=90,width=\textwidth,height=\textheight,keepaspectratio]{#1}}%
    \vfil}%
  \newpage}

% ---------------------------------------------------------------------------
% Chapitres (division éditoriale)
% ---------------------------------------------------------------------------
\newcommand{\LJchapitre}[3]{%
  \LJpageimpaire
  \renewcommand{\LJmois}{#3}\renewcommand{\LJfolio}{Chapitre~#1}%
  \chapter{#2}%
  \begin{center}\small\itshape #3\end{center}\vspace{1mm}}
% Ouverture de chapitre (page paire, en regard du début de l'introduction) : #1 numéro ; #2 titre ; #3 période ;
% #4 carte (PDF) ; #5 illustrations sous la carte
\newcommand{\LJouverture}[5]{%
  \LJpagepaire
  \renewcommand{\LJmois}{#3}\renewcommand{\LJfolio}{Chapitre~#1}%
  \thispagestyle{plain}%
  \refstepcounter{chapter}\phantomsection
  \addcontentsline{toc}{chapter}{\protect\numberline{\thechapter}#2}%
  \begin{center}
    {\small\scshape chapitre~#1}\par\vspace{0.8mm}
    {\Large\scshape #2}\par\vspace{0.8mm}
    {\small\itshape #3}\par
  \end{center}\vspace{-1mm}
  \noindent\includegraphics[width=\textwidth]{#4}\par
  \vfill
  \noindent #5\par
  \clearpage}
% Photographie légendée (dans une ouverture ou une page d'illustration) ; #1 largeur ; #2 fichier ; #3 légende
\newcommand{\LJphoto}[3]{%
  \begin{minipage}[b]{#1}\centering
    \includegraphics[width=\linewidth,height=44mm,keepaspectratio]{#2}\par\vspace{0.8mm}
    {\scriptsize\raggedright #3\par}
  \end{minipage}}
% Fin d'une introduction : si la page suivante est impaire (la double page de transcription ne pourrait pas
% commencer), elle reçoit l'illustration #1 au lieu de rester blanche.
\newcommand{\LJfinintro}[1]{\clearpage\ifodd\value{page}\thispagestyle{plain}\vspace*{\fill}#1\vspace*{\fill}\clearpage\fi}

% Transcription d'une lettre en annexe (paragraphes de l'original ; police de la transcription)
\newenvironment{LJlettre}{\par\vspace{1mm}\begingroup\LJpolice\small
  \setlength{\leftskip}{4mm}\setlength{\rightskip}{2mm plus 1fil}\setlength{\parindent}{1.2em}%
  \setlength{\parskip}{0.4mm}\emergencystretch=2em}{\par\endgroup\vspace{1mm}}

% ---------------------------------------------------------------------------
% Index
% ---------------------------------------------------------------------------
\newcommand{\LJindex}[2]{%
  \clearpage\renewcommand{\LJmois}{Index}\renewcommand{\LJfolio}{#1}%
  \chapter*{#1}\addcontentsline{toc}{chapter}{#1}%
  {\footnotesize\setlength{\parindent}{0pt}\setlength{\parskip}{0.6pt}%
   \begin{multicols}{2}\raggedright #2\end{multicols}}}
\newcommand{\LJentree}[2]{\hangindent=1em\hangafter=1 #1 : {\LJchiffres #2}\par}

% ---------------------------------------------------------------------------
% Table de concordance des feuillets et des clichés
% ---------------------------------------------------------------------------
\newcommand{\LJtableconcordance}[2]{%
  \clearpage\renewcommand{\LJmois}{Annexe}\renewcommand{\LJfolio}{Concordance des feuillets et des clichés}%
  \chapter*{Concordance des feuillets et des clichés}\addcontentsline{toc}{chapter}{Concordance des feuillets et des clichés}%
  {\footnotesize Feuillet et face restitués dans l'ordre de lecture ; numéro du cliché ;
  date(s) de la page ; page de la présente édition. \textit{inv.} : feuillet photographié verso d'abord ;
  \textit{trav.} : page écrite en travers ; \textit{seule face} : l'autre face du feuillet n'a pas été photographiée. #2\par}
  \vspace{2mm}
  {\fontsize{7.4pt}{8.8pt}\selectfont\LJchiffres
  \setlength{\tabcolsep}{4pt}
  \begin{longtable}{@{}l r p{72mm} r l@{}}
  \toprule
  \textbf{Feuillet} & \textbf{Cliché} & \textbf{Date} & \textbf{Page} & \\
  \midrule\endhead
  \bottomrule\endfoot
  #1
  \end{longtable}}}
